\section{General background}\label{sec-diabat-common}
Adiabatic-to-diabatic (ATD) diabatization constructs quasi-diabatic states from selected adiabatic states of a molecular system. If \(\lvert\Psi_n\rangle\) are the chosen adiabatic wave functions, the resulting states are
\begin{equation}
\lvert\Phi_k\rangle=\sum_n U_{nk}\lvert\Psi_n\rangle,
\end{equation}
where $\hat{U}$ is the ATD transformation. In the adiabatic representation, the Hamiltonian is diagonal when the original adiabatic energies are used. The quasi-diabatic Hamiltonian $\hat{H}^{\mathrm{d}}$ is obtained via
\begin{equation}
\hat{H}^{\rm d}=\hat{U}^{\dagger}\hat{H}^{\rm ad}\hat{U}.
\end{equation}
Its diagonal elements are quasi-diabatic state energies and its off-diagonal elements are electronic couplings.

\progname{} implements FPHD, GMH, FCD, FED, and multistate FED--FCD. Each method selects $\hat{U}$ and idenfity the state character using a different state descriptor. Their scope and restrictions are given in the following sections. For multistate calculations, different states can have similar descriptors. The program therefore normally performs a partial Hamiltonian diagonalization within groups of states having the same assigned character. The resulting transformation and Hamiltonian should be interpreted together.
See the software release paper\cite{Diabat2026} for further theoretical details.

The program uses a common phase convention for coupling signs at a fixed nuclear geometry and Cartesian orientation.
Although this convention often leads to consistent phases between nearby geometries, it does not itself guarantee phase continuity when the nuclear geometry changes.
 Appendix~\ref{app-phase} explains the convention.
 The Phasefix utility is described in Section~\ref{sec-phasefix}.